\documentclass[11pt]{article}
\usepackage[utf8]{inputenc}
\usepackage[english]{babel}
\usepackage{amsmath}
\usepackage{graphicx}
\usepackage{float}
\usepackage{lipsum}
\usepackage{multicol}
\usepackage{xcolor}
\usepackage{tabularx}
\usepackage{booktabs}
\usepackage{hyperref}
\newcolumntype{Y}{>{\centering\arraybackslash}X}
\usepackage[left=2.00cm, right=2.00cm, top=2.00cm, bottom=2.00cm]{geometry}
\usepackage{enumitem}

\title{AN2DL Reports Template}

\begin{document}

\begin{figure}[H]
      \raggedright
      \includegraphics[scale=0.4]{polimi.png} \hfill
      \includegraphics[scale=0.3]{airlab.jpeg}
\end{figure}

\vspace{5mm}

\begin{center}
      % Select between First and Second
      {\Large \textbf{AN2DL - Second Homework Report}}\\
      \vspace{2mm}
      % Change with your Team Name
      {\Large \textbf{LosPollosHermanos}}\\
      \vspace{2mm}
      % Team Members Information
      {\large Mohammadhossein Allahakbari,}
      {\large Michele Miotti,}
      {\large Francesco Pesce}\\
      \vspace{2mm}
      % Codabench Nicknames
      {mh2033,}
      {michelem,}
      {francescopesce}\\
      \vspace{2mm}
      % Matriculation Numbers
      {246639,}
      {249499,}
      {247974}\\
      \vspace{5mm}
      \today
\end{center}
\vspace{5mm}

\begin{multicols}{2}
      % Note: The following sections represent a suggested
      % structure. We don't need to follow it strictly.

      % -----------------------------------------------------------------------
      % INTRODUCTION
      % -----------------------------------------------------------------------
      \section{Introduction}
      % In this section, you should present your project's context and
      % objectives. You might want to:
      % \begin{itemize}
      %     \item Dehe problem (\textit{you may use italics to highlight
      %               definitions})
      %     \item State your goals (\textbf{emphasise key points with bold})
      %     \item Outline your approach
      % \end{itemize}

      % \noindent For instance, you might write: ``This project focuses on
      % \textit{image classification} using \textbf{deep learning} techniques."

      This report presents the results of the \textit{semantic segmentation}
      task proposed in the second homework of the \texttt{Artificial Neural Networks and Deep Learning} course. The goal of the project is to categorize all the pixels in a 64x128 grayscale image into different classes, each representing a different type of terrain.

      % -----------------------------------------------------------------------
      % PROBLEM ANALYSIS
      % -----------------------------------------------------------------------
      \section{Problem Analysis}
      % Here you can discuss your initial analysis of the problem. Consider
      % including:
      % \begin{enumerate}
      %     % 8 classes, 96x96 rgb images, labels, etc
      %     \item Dataset characteristics
      %     \item Main challenges % The test set is horrible
      %     That the test set was not horrible? Is this what they mean?
      %     \item Initial assumptions 
      % \end{enumerate}

      % \noindent If you need to reference papers, use the citation command:
      % Recent work~\cite{lecun2015deep} suggests..."

      The dataset is composed of 64x128 grayscale images, each pixel belonging to one of five classes (background, soil, bedrock, sand, big rock). Without considering the background, the model is evaluated on the intersection over union (IoU) metric, which is the ratio between the intersection and the union of the predicted and true masks. The minimum IoU is chosen within the classes as the final score.

      Labels are only provided in the training set, and the test set (available on Kaggle) is only used for the evaluation of the model. In order to validate the model, we split the training set into a training and validation set, using a 80-20 ratio. Differently from the first homework, we noticed that the validation and test sets' performances were similar.

      After manually analyzing the dataset, we found that a portion of the test presented 110 irrelevant, or otherwise unnecessary images. We didn't consider said images in our training, as we didn't expect them to be present in the test set. Moreover, the provided labels do not appear to be perfect, as they were labelled by hand. This could lead to a decrease in the model's performance which is not due to the model itself.

      % -----------------------------------------------------------------------   
      % METHOD
      % -----------------------------------------------------------------------
      \section{Method}
      % Not sure what they are asking here. The final model?
      % This section should detail your approach. You can use equations to
      % explain your methodology. For example, a simple model representation:
      % \begin{equation}
      %     \label{eq:model}
      %     f(x) = \text{softmax}(Wx + b)
      % \end{equation}

      % \noindent Or a more complex loss function:
      % \begin{equation}
      %     \label{eq:loss}
      %     \mathcal{L} = -\frac{1}{N}\sum_{i=1}^{N} y_i\log(\hat{y}_i)
      % \end{equation}

      % \noindent Reference these equations in your text, like:``As shown in
      % equation~\ref{eq:model}..."

      Most of the development process was done on a \texttt{Jupyter} notebook, using the \texttt{Tensorflow}\cite{TensorFlow} and \texttt{Keras}\cite{chollet2015keras} libraries for \texttt{Python}, mostly run locally using an \texttt{NVIDIA RTX 2060 Mobile} for training, due to strict GPU usage limits of popular platforms, such as \texttt{Google Colab}.

      % -----------------------------------------------------------------------
      % EXPERIMENTS
      % -----------------------------------------------------------------------
      \label{sec:experiments}
      \section{Experiments}
      We performed many experiments, but most of them have the same basic structure, which was heavily inspired by the U-Net architecture. All models have a "V" shape, with a contracting path on the left and an expanding path on the right. The contracting path is composed of a series of convolutional layers, where each layer halves the spatial dimensions of the input, while doubling the number of channels. The expanding path is composed of a series of up-sampling layers, which act as the opposite of the previous layers, doubling the spatial dimensions and halving the number of channels.

      The latest models also feature connections between the contracting and expanding paths of the same size, which are used to pass information from the contracting to the expanding section. In chronological order, the main breakthroughs were:
      \begin{itemize}[leftmargin=*]
            \setlength\itemsep{0em}
            \item \textbf{Experiment:} Description
      \end{itemize}

      Some experiments did not lead to any improvement, and were discarded. For example, ...

      % For your experiments, you might want to present your results in tables.
      % Here's an example of a wide table comparing different models:

      % \begin{table*}[t]
      %     \centering
      %     \setlength{\tabcolsep}{3pt}
      %     \caption{An example of wide table. Best results are highlighted in
      %         \textbf{bold}.}
      %     \begin{tabularx}{\textwidth}{lYYYc}
      %         \toprule
      %         Model            & Accuracy                  & Precision
      %                          & Recall                    & ROC AUC
      %         \\
      %         \midrule
      %         VGG18            & 72.20 $\pm$ 3.06          & 94.95 $\pm$ 0.52
      %                          &
      %         86.95 $\pm$ 0.55 & 80.16 $\pm$ 0.81
      %         \\
      %         Custom Model     & 27.71 $\pm$ 3.19          & 75.70 $\pm$ 1.07
      %                          & 55.75 $\pm$ 2.16          & 36.60 $\pm$ 1.26
      %         \\
      %         ResNet18         & \textbf{89.24 $\pm$ 2.38} & \textbf{95.54
      %         $\pm$ 0.49}      & \textbf{93.43 $\pm$ 1.30} & \textbf{91.68 $\pm$
      %             0.71}
      %         \\
      %         \bottomrule
      %     \end{tabularx}
      %     \label{tab:Performance}
      % \end{table*}

      % \noindent For more specific measurements, you might use a narrower
      % table:

      % \begin{table}[H]
      %     \centering
      %     \setlength{\tabcolsep}{3pt}
      %     \caption{An example of table. Best results may be highlighted in
      %         \textbf{bold}.}
      %     \begin{tabularx}{\linewidth}{lY}
      %         \toprule
      %         Time [$\mu$s] & Distance [mm] \\
      %         \midrule
      %         22$\pm$4      & 8$\pm$1       \\
      %         17$\pm$3      & 7$\pm$1       \\
      %         15$\pm$3      & 6$\pm$1       \\
      %         13$\pm$2      & 5$\pm$1       \\
      %         10$\pm$2      & 4$\pm$1       \\
      %         8$\pm$2       & 3$\pm$1       \\
      %         5$\pm$1       & 2$\pm$1       \\
      %         37$\pm$1      & 1$\pm$1       \\
      %         \bottomrule
      %     \end{tabularx}
      %     \label{tb:Measurements}
      % \end{table}

      % \noindent You can also include figures to visualise your results:
      % \begin{figure}[H]
      %     \centering
      %     \includegraphics[width=0.75\linewidth]{random.jpeg}
      %     \caption{Example figure showing [describe what the figure shows]}
      %     \label{fig:results}
      % \end{figure}

      % \noindent Reference figures using like:``As shown in
      % Figure~\ref{fig:results}..."

      % -----------------------------------------------------------------------
      % RESULTS
      % -----------------------------------------------------------------------
      \label{sec:results}
      \section{Results}
      % Present your main findings here. You might want to:
      % \begin{itemize}
      %     \item Compare your results with baselines
      %     \item Highlight key achievements using \textbf{bold text}
      %     \item Explain any unexpected outcomes
      % \end{itemize}

      % -----------------------------------------------------------------------
      % DISCUSSION
      % -----------------------------------------------------------------------
      \section{Discussion}
      % In this section, analyse your results critically. Consider:
      % \begin{itemize}
      %     \item Strengths and weaknesses
      %     \item Limitations and assumptions
      % \end{itemize}

      % -----------------------------------------------------------------------
      % CONTRIBUTIONS
      % -----------------------------------------------------------------------
      \section{Contributions}

      Most ideas and propositions were evenly distributed among the team members, while actual coding was more specific to each member's strengths. Many different branches were created in the project's repository to experiment with different ideas and techniques, and the final model was a result of combining the best experiments from each branch.

      % -----------------------------------------------------------------------
      % CONCLUSIONS
      % -----------------------------------------------------------------------
      \section{Conclusions}
      % Summarise your work and discuss potential future directions. This is
      % where you can:
      % \begin{itemize}
      %     \item Restate main contributions
      %     \item Suggest improvements
      %     \item Propose future work
      % \end{itemize}

      % Remember to include the bibliography!
      \bibliography{references}
      \bibliographystyle{abbrv}

\end{multicols}
\end{document}